%!TeX root=../tese.tex
%("dica" para o editor de texto: este arquivo é parte de um documento maior)
% para saber mais: https://tex.stackexchange.com/q/78101

% As palavras-chave são obrigatórias, em português e em inglês, e devem ser
% definidas antes do resumo/abstract. Acrescente quantas forem necessárias.
\palavraschave{Biblioteca, Framework, Reutilização de código, Ciclo de vida, Popularidade, Mineração de repositórios}

\keywords{Library, Framework, Code reuse, Lifecycle, Popularity, Repository mining}

% O resumo é obrigatório, em português e inglês. Estes comandos também
% geram automaticamente a referência para o próprio documento, conforme
% as normas sugeridas da USP.
\resumo{
O desenvolvimento de software tem como pilar o uso de bibliotecas e frameworks de código aberto, que permitem a reutilização de código e aceleram o processo de desenvolvimento, mas a literatura ainda carece de uma metodologia que avalie, de maneira unificada, tanto o ciclo de vida quanto a popularidade desse tipo de ferramenta. A proposta desse trabalho é apresentar uma metodologia capaz de avaliar de maneira unificada e sistemática métricas de ciclo de vida e de popularidade de bibliotecas e frameworks de código aberto, com o objetivo de identificar tendências de adoção ou abandono de ferramentas pela comunidade. Esse trabalho está dividido em quatro etapas: (I) Minerar repositórios para levantar métricas de ciclo de vida e popularidade; (II) Proposta de um método de avaliação utilizando modelos temporais; (III) Comparação dos resultados com resultados de outros métodos existentes na literatura; e (IV) Validação do método em cenários reais de desenvolvimento. Com isso, objetiva-se fornecer uma metodologia que auxilie a identificação de tendências de adoção ou abandono de bibliotecas e frameworks de código aberto, contribuindo para a melhoria na tomada de decisão sobre qual ferramenta utilizar.
}

\abstract{
Software development relies heavily on open-source libraries and frameworks, which enable code reuse and accelerate the development process, however, the literature still lacks a methodology that evaluates both the lifecycle and the popularity of such tools in a unified manner. This work proposes a methodology capable of evaluating lifecycle and popularity metrics for open-source libraries and frameworks in a unified and systematic way, aiming to identify community trends regarding tool adoption or abandonment. The work is divided into four stages: (I) Mining repositories to gather lifecycle and popularity metrics; (II) Proposing an evaluation method using temporal models; (III) Comparing the results with those from other methods found in the literature; and (IV) Validating the method in real-world development scenarios. The goal is to provide a methodology that assists in identifying trends in the adoption or abandonment of open-source libraries and frameworks, thereby contributing to better decision-making regarding which tools to use.
}
